\chapter{Theoretischer Rahmen: \DO als soziotechnisches Phänomen}
\label{ch:theoretical-framework}

Ein fundiertes theoretisches Verständnis von \DO als soziotechnisches System ist unerlässlich, um die nachfolgend analysierten kulturellen und sozialen Faktoren korrekt einordnen zu können.
Eine rein technische Betrachtung greift zu kurz und verkennt die fundamentalen organisatorischen und menschlichen Dynamiken, die den Erfolg einer \DO-Transformation bestimmen.
Dieses Kapitel verortet \DO daher in seinem historischen und konzeptionellen Kontext und beleuchtet es aus organisationssoziologischer Perspektive.

\section{Ursprung und Prinzipien von \DO}

Der Ursprung von \DO liegt in einer kleinen Gruppe von Softwareentwicklern rund um Patrick Debois.
Sie fanden sich 2009 zusammen, um die traditionellen Barrieren zwischen Entwicklung und Betrieb zu überwinden, und legten damit den Grundstein für die \DO-Bewegung~\cite[10]{jha23}.

Die Kernprinzipien von \DO werden oft in einem konzeptionellen Modell zusammengefasst, das als Akronym CALMS bekannt ist. Diese stehen für:

\begin{itemize}
	\item \textbf{C}ulture (Kultur): Die Förderung von Zusammenarbeit, gemeinsamer Verantwortung und Vertrauen.
	\item \textbf{A}utomation (Automatisierung): Die Automatisierung von repetitiven Aufgaben in der Software"=Lieferkette.
	\item \textbf{L}ean (Schlanke Prozesse): Die Anwendung von Lean"=Prinzipien zur Eliminierung von Verschwendung und zur Optimierung von Wertströmen.
	\item \textbf{M}easurement (Messung): Die kontinuierliche Messung von Prozessen und Systemen, um datengestützte Entscheidungen zu ermöglichen.
	\item \textbf{S}haring (Teilen): Der offene Austausch von Wissen, Werkzeugen und Erfahrungen über Team"= und Abteilungsgrenzen hinweg.
\end{itemize}
Dieses Modell versucht, eine Balance zwischen sozialen Dimensionen (Kultur, Teilen) und technischen Dimensionen (Automatisierung) herzustellen und betont, dass die Kultur die Basis für alle weiteren Prinzipien bildet~\cite[14\psq]{jha23}.

\section{Soziotechnische Systemtheorie}

\DO ist ein klassisches Beispiel für ein soziotechnisches System nach \textcite[135-150]{rop09}.
Die soziotechnische Systemtheorie besagt, dass soziale und technische Subsysteme in einer Organisation untrennbar miteinander verbunden sind und daher auch bei Transformationen gemeinsam angepasst werden müssen, ihr volles Potenzial auszuschöpfen.
Im Kontext von \DO lässt sich eine Organisation in zwei eng verflochtene Subsysteme unterteilen:

Das \emph{soziale Subsystem} beinhaltet einerseits die Menschen, Teams, Kommunikationsstrukturen, Führungsstile und die Organisationskultur.
Das \emph{technische Subsystem} umfasst andererseits die Werkzeuge, Automatisierungspipelines, Infrastruktur"=as"=Code und Monitoring"=Systeme.

Der entscheidende Fehler bei vielen \DO-Initiativen ist die einseitige Fokussierung auf das technische Subsystem.
Praktiker warnen eindringlich davor, \DO fälschlicherweise ausschließlich auf Werkzeuge und Automatisierung zu reduzieren~\cite[8]{aza23}.
Ein solcher Ansatz ignoriert die sozialen Dynamiken, Widerstände und kulturellen Voraussetzungen, die letztlich über Erfolg oder Misserfolg entscheiden.

\section{Komponenten der Organisationskultur}
Der Erfolg von \DO hängt also weniger von der technischen Umsetzung ab, als von der Etablierung einer Kultur, die Silos aufbricht und gemeinsame Ziele in den Vordergrund stellt~\cites[7\psq]{aza23}[313]{nay24}.
Silos entstehen häufig durch unterschiedliche Abteilungsziele, starre Strukturen, Kommunikationsbarrieren oder gar Misstrauen zwischen Teams, sodass diese getrennt voneinander arbeiten und damit die für \DO essenzielle Zusammenarbeit behindern~\cite[5]{mak25}.
Die Literatur identifiziert dabei konsistent eine Reihe von Kernattributen, die eine \DO-Kultur auszeichnen~\cites[4]{kre22}[281]{mas18}:

\begin{itemize}
	\item \emph{Kommunikation und Kollaboration:} Der offene und ständige Austausch zwischen allen Beteiligten.
	\item \emph{Kontinuierliches Feedback:} Schnelle und regelmäßige Rückmeldungen auf allen Ebenen.
	\item \emph{Geteilte Verantwortung:} Teams übernehmen gemeinsam die Verantwortung für den gesamten Lebenszyklus einer Anwendung~\cite[284\psq]{kar24}.
	\item \emph{Vertrauen:} Ein hohes Maß an Vertrauen zwischen den Teams und Individuen als Grundlage für effektive Zusammenarbeit~\cite[8]{aza23}.
\end{itemize}
Diese Attribute sind nicht isoliert zu betrachten, sondern bilden das Gerüst für das übergeordnete Ziel: den Aufbau einer \emph{kollaborativen Kultur}, in der Teams nicht gegeneinander, sondern miteinander für den Erfolg des Produkts arbeiten.

\paragraph{}
Nachdem das theoretische Fundament gelegt ist, das \DO als ein von Kultur und sozialen Interaktionen geprägtes System begreift, widmet sich das folgende Kapitel den konkreten Faktoren, die sich in der Praxis als entscheidend für den Erfolg erwiesen haben.